\begingroup
\renewcommand{\arraystretch}{1.45}
\fancyhead[C]{}
\providecommand{\TStart}{\hspace*{.5em}}\providecommand{\TCont}{\hspace*{1em}}\providecommand{\TD}{\hspace{.3em}.\hspace{.3em}.}
\clearpage\fancyhead[L]{\small CONTENTS.}\pdfbookmark[0]{Contents}{original-inhalt}\SourceLeaf{506}


\begin{center}
\textbf{CONTENTS.}\par
\medskip
\textbf{GAUSS’S WORKS VOLUME III.\quad ANALYSIS.}\par
\medskip
\rule{2em}{0.4pt}
\end{center}
\begingroup\small
\setlength{\tabcolsep}{0pt}
\setlength{\LTleft}{0pt}\setlength{\LTright}{0pt}
\begin{longtable}{@{}l@{\hspace{.01\linewidth}}p{.19\linewidth}@{\hspace{.01\linewidth}}p{.075\linewidth}@{\hspace{.005\linewidth}}p{.05\linewidth}@{}}
\multicolumn{4}{c}{\textbf{ALGEBRAIC FUNCTIONS.}}\\*[.3em]
\multicolumn{4}{@{}l@{}}{\textit{Treatises.}}\\*
\parbox[b]{.64\linewidth}{\raggedright New Proof of the Theorem That Every Rational Integral Algebraic Function of One Variable Can Be Resolved into Real Factors of the First or Second Degree ~\hfill\mbox{.\kern.08em .\kern.08em .\kern.08em .\kern.08em .\kern.08em .\kern.08em .\kern.08em .\kern.08em .\kern.08em .\kern.08em .\kern.08em .\kern.08em .\kern.08em .\kern.08em .\kern.08em .\kern.08em .\kern.08em .\kern.08em .\kern.08em .\kern.08em .}\par} & 1799 & Page & 1\\
\parbox[b]{.64\linewidth}{\raggedright Second New Proof of the Theorem That Every Rational Integral Algebraic Function of One Variable Can Be Resolved into Real Factors of the First or Second Degree ~\hfill\mbox{.\kern.08em .\kern.08em .\kern.08em .\kern.08em .\kern.08em .\kern.08em .\kern.08em .\kern.08em .\kern.08em .\kern.08em .\kern.08em .\kern.08em .\kern.08em .\kern.08em .\kern.08em .\kern.08em .\kern.08em .\kern.08em .\kern.08em .}\par} & 1815 Dec.\TD & --- & 31\\
\parbox[b]{.64\linewidth}{\raggedright Third Proof of the Theorem on the Resolution of Integral Algebraic Functions into Real Factors. Supplement to the Preceding Treatise\par} & 1816 Jan.\TD & --- & 57\\
\parbox[b]{.64\linewidth}{\raggedright Proof of an Algebraic Theorem ~\hfill\mbox{.\kern.08em .\kern.08em .\kern.08em .\kern.08em .\kern.08em .\kern.08em .\kern.08em .\kern.08em .\kern.08em .\kern.08em .\kern.08em .\kern.08em .}\par} & 1828 & --- & 65\\
\parbox[b]{.64\linewidth}{\raggedright Contributions to the Theory of Algebraic Equations ~\hfill\mbox{.\kern.08em .\kern.08em .\kern.08em .\kern.08em .\kern.08em .\kern.08em .\kern.08em .}\par} & 1849 July\TD & --- & 71\\
\multicolumn{4}{@{}l@{}}{\textit{Notices of His Own Treatises.}}\\*
\parbox[b]{.64\linewidth}{\raggedright Second New Proof of the Theorem That Every Rational Integral Algebraic Function of One Variable Can Be Resolved into Real Factors of the First or Second Degree ~\hfill\mbox{.\kern.08em .\kern.08em .\kern.08em .\kern.08em .\kern.08em .\kern.08em .\kern.08em .\kern.08em .\kern.08em .\kern.08em .\kern.08em .\kern.08em .\kern.08em .\kern.08em .\kern.08em .\kern.08em .\kern.08em .\kern.08em .}\par} & 1815 Dec.\TD & --- & 105\\
\parbox[b]{.64\linewidth}{\raggedright Third Proof of the Theorem on the Resolution of Integral Algebraic Functions into Real Factors ~\hfill\mbox{.\kern.08em .\kern.08em .\kern.08em .\kern.08em .\kern.08em .\kern.08em .\kern.08em .\kern.08em .\kern.08em .\kern.08em .\kern.08em .\kern.08em .\kern.08em .}\par} & 1816 March\TD & --- & 107\\
\parbox[b]{.64\linewidth}{\raggedright Contributions to the Theory of Algebraic Equations ~\hfill\mbox{.\kern.08em .\kern.08em .\kern.08em .\kern.08em .\kern.08em .\kern.08em .\kern.08em .}\par} & 1849 July\TD & --- & 113\\
\multicolumn{4}{@{}l@{}}{\textit{Notices of Writings by Others.}}\\*
\parbox[b]{.64\linewidth}{\raggedright \textsc{Coytier.} Mathematical Investigations ~\hfill\mbox{.\kern.08em .\kern.08em .\kern.08em .\kern.08em .\kern.08em .\kern.08em .\kern.08em .\kern.08em .\kern.08em .\kern.08em .\kern.08em .\kern.08em .}\par} & 1813 March\TD & --- & 116\\
\parbox[b]{.64\linewidth}{\raggedright \textsc{Fourier.} Analysis of Determinate Equations ~\hfill\mbox{.\kern.08em .\kern.08em .\kern.08em .\kern.08em .\kern.08em .\kern.08em .\kern.08em .\kern.08em .\kern.08em .\kern.08em .}\par} & 1833 Febr.\TD & --- & 119\\[1em]
\multicolumn{4}{c}{\textbf{GAUSS’S SERIES.}}\\*
\multicolumn{4}{@{}l@{}}{\textit{Treatises.}}\\*
\multicolumn{4}{@{}l@{}}{\TStart General Investigations Concerning the Infinite Series}\\[.3em]
\noalign{\nobreak}
\multicolumn{4}{@{}l@{}}{\TCont $\displaystyle\begin{gathered}
1+\frac{\alpha.\beta}{1.\gamma}x
+\frac{\alpha(\alpha+1)\beta(\beta+1)}{1.2.\gamma(\gamma+1)}xx
\\[9pt]+\frac{\alpha(\alpha+1)(\alpha+2)\beta(\beta+1)(\beta+2)}{1.2.3.\gamma(\gamma+1)(\gamma+2)}x^{3}
+\mathord{.}\,\mathord{.}\end{gathered}$}\\[.3em]
\noalign{\nobreak}
\parbox[b]{.64\linewidth}{\raggedright First Part ~\hfill\mbox{.\kern.08em .\kern.08em .\kern.08em .\kern.08em .\kern.08em .\kern.08em .\kern.08em .\kern.08em .\kern.08em .\kern.08em .\kern.08em .\kern.08em .\kern.08em .\kern.08em .\kern.08em .\kern.08em .\kern.08em .\kern.08em .\kern.08em .\kern.08em .\kern.08em .}\par} & 1812 Jan.\TD & --- & 123\\
\end{longtable}
\endgroup


\SourceLeaf{507}


\begingroup\small
\setlength{\tabcolsep}{0pt}
\setlength{\LTleft}{0pt}\setlength{\LTright}{0pt}
\begin{longtable}{@{}l@{\hspace{.01\linewidth}}p{.19\linewidth}@{\hspace{.01\linewidth}}p{.075\linewidth}@{\hspace{.005\linewidth}}p{.05\linewidth}@{}}
\parbox[b]{.64\linewidth}{\raggedright New Method of Finding the Values of Integrals by Approximation ~\hfill\mbox{.\kern.08em .\kern.08em .}\par} & 1814 Sept.\TD & Page & 163\\
\multicolumn{4}{@{}l@{}}{\textit{Notices of His Own Treatises.}}\\*
\parbox[b]{.64\linewidth}{\raggedright General Investigations Concerning the Infinite Series etc. First Part. ~\hfill\mbox{.\kern.08em .\kern.08em .\kern.08em .\kern.08em .}\par} & 1812 Febr.\TD & --- & 197\\
\parbox[b]{.64\linewidth}{\raggedright New Method of Finding the Values of Integrals by Approximation ~\hfill\mbox{.\kern.08em .\kern.08em .}\par} & 1814 Sept.\TD & --- & 202\\
\multicolumn{4}{@{}l@{}}{\textit{Manuscript Estate.}}\\*
\parbox[b]{.64\linewidth}{\raggedright Determination of Our Series by a Differential Equation of the Second Order ~\hfill\mbox{.}\par} & & --- & 207\\
\parbox[b]{.64\linewidth}{\raggedright Remarks on This Treatise ~\hfill\mbox{.\kern.08em .\kern.08em .\kern.08em .\kern.08em .\kern.08em .\kern.08em .\kern.08em .\kern.08em .\kern.08em .\kern.08em .\kern.08em .\kern.08em .\kern.08em .}\par} & & --- & 230\\[1em]
\multicolumn{4}{@{}p{.98\linewidth}@{}}{\centering\textbf{COMMUNICATIONS ON VARIOUS WRITINGS.}}\\*[.4em]
\parbox[b]{.64\linewidth}{\raggedright \textsc{Pfaff.} General Method for Completely Integrating Partial Differential Equations and Ordinary Differential Equations, Both of the First Order, among Any Number of Variables ~\hfill\mbox{.\kern.08em .\kern.08em .\kern.08em .\kern.08em .\kern.08em .\kern.08em .\kern.08em .\kern.08em .\kern.08em .\kern.08em .\kern.08em .}\par} & 1815 July\TD & --- & 231\\
\parbox[b]{.64\linewidth}{\raggedright \textsc{Callet.} Logarithm Tables ~\hfill\mbox{.\kern.08em .\kern.08em .\kern.08em .\kern.08em .\kern.08em .\kern.08em .\kern.08em .\kern.08em .\kern.08em .\kern.08em .\kern.08em .\kern.08em .\kern.08em .\kern.08em .\kern.08em .}\par} & 1802 Nov.\TD & --- & 241\\
\parbox[b]{.64\linewidth}{\raggedright \textsc{Prasse.} Logarithm Tables for Numbers, Sines and Tangents ~\hfill\mbox{.\kern.08em .}\par} & 1811 May\TD & --- & 241\\
\parbox[b]{.64\linewidth}{\raggedright \textsc{Prasse.} Logarithm Tables for Numbers, Sines and Tangents\par} & 1814 Dec.\TD & --- & 243\\
\parbox[b]{.64\linewidth}{\raggedright \textsc{Gauss.} Tables for More Conveniently Calculating the Logarithm of the Sum or Difference of Two Quantities That Are Themselves Given Only by Their Logarithms ~\hfill\mbox{.\kern.08em .\kern.08em .\kern.08em .\kern.08em .\kern.08em .\kern.08em .\kern.08em .\kern.08em .\kern.08em .\kern.08em .\kern.08em .\kern.08em .\kern.08em .\kern.08em .\kern.08em .\kern.08em .\kern.08em .\kern.08em .\kern.08em .\kern.08em .\kern.08em .}\par} & 1812 Nov.\TD & --- & 244\\
\parbox[b]{.64\linewidth}{\raggedright \textsc{Pasquich.} Abridged Logarithmic and Trigonometric Tables ~\hfill\mbox{.\kern.08em .\kern.08em .\kern.08em .\kern.08em .}\par} & 1817 Oct.\TD & --- & 246\\
\parbox[b]{.64\linewidth}{\raggedright \relax[\textsc{Matthiessen.}] Table for More Conveniently Calculating the Logarithm of the Sum or Difference of Two Quantities That Are Themselves Given Only by Their Logarithms ~\hfill\mbox{.\kern.08em .\kern.08em .\kern.08em .\kern.08em .\kern.08em .\kern.08em .\kern.08em .\kern.08em .\kern.08em .\kern.08em .\kern.08em .\kern.08em .\kern.08em .\kern.08em .\kern.08em .\kern.08em .\kern.08em .\kern.08em .}\par} & 1819 Jan.\TD & --- & 250\\
\parbox[b]{.64\linewidth}{\raggedright Logarithm Tables with Six Digits ~\hfill\mbox{.\kern.08em .\kern.08em .\kern.08em .\kern.08em .\kern.08em .\kern.08em .\kern.08em .\kern.08em .\kern.08em .\kern.08em .\kern.08em .\kern.08em .\kern.08em .\kern.08em .}\par} & 1822 Dec.\TD & --- & 252\\
\parbox[b]{.64\linewidth}{\raggedright \textsc{Babbage.} Table of logarithms of the natural numbers ~\hfill\mbox{.\kern.08em .\kern.08em .\kern.08em .\kern.08em .\kern.08em .\kern.08em .\kern.08em .}\par} & 1828 Jan.\TD & --- & 253\\
\parbox[b]{.64\linewidth}{\raggedright \textsc{Hassler.} Logarithmic and Trigonometric Tables ~\hfill\mbox{.\kern.08em .\kern.08em .\kern.08em .\kern.08em .\kern.08em .\kern.08em .\kern.08em .\kern.08em .}\par} & 1831 March\TD & --- & 255\\
\parbox[b]{.64\linewidth}{\raggedright \textsc{Vega} and \textsc{Hülsse.} Collection of Mathematical Tables ~\hfill\mbox{.\kern.08em .\kern.08em .\kern.08em .\kern.08em .\kern.08em .}\par} & 1840 & --- & 255\\
\parbox[b]{.64\linewidth}{\raggedright \textsc{Vega.} Treasury of Logarithms ~\hfill\mbox{.\kern.08em .\kern.08em .\kern.08em .\kern.08em .\kern.08em .\kern.08em .\kern.08em .\kern.08em .\kern.08em .\kern.08em .\kern.08em .\kern.08em .\kern.08em .\kern.08em .}\par} & 1851 May\TD & --- & 257\\[1em]
\multicolumn{4}{c}{\textbf{INTERPOLATION.}}\\*
\multicolumn{4}{@{}l@{}}{\textit{Manuscript Estate.}}\\*
\parbox[b]{.64\linewidth}{\raggedright Theory of Interpolation Treated by a New Method ~\hfill\mbox{.\kern.08em .\kern.08em .\kern.08em .\kern.08em .\kern.08em .\kern.08em .\kern.08em .\kern.08em .\kern.08em .\kern.08em .\kern.08em .\kern.08em .\kern.08em .\kern.08em .\kern.08em .}\par} & & --- & 265\\
\parbox[b]{.64\linewidth}{\raggedright Remarks on This Treatise ~\hfill\mbox{.\kern.08em .\kern.08em .\kern.08em .\kern.08em .\kern.08em .\kern.08em .\kern.08em .\kern.08em .\kern.08em .\kern.08em .\kern.08em .\kern.08em .\kern.08em .\kern.08em .\kern.08em .\kern.08em .\kern.08em .\kern.08em .}\par} & & --- & 328\\[1em]
\multicolumn{4}{c}{\textbf{ELLIPTIC FUNCTIONS.}}\\*
\multicolumn{4}{@{}l@{}}{\textit{Treatise.}}\\*
\parbox[b]{.64\linewidth}{\raggedright Determination of the Attraction,. That a Planet Would Exert on Any Given Point If Its Mass Were Uniformly Distributed over Its Whole Orbit in Proportion to the Time in Which the Individual Parts Are Described ~\hfill\mbox{.\kern.08em .\kern.08em .\kern.08em .\kern.08em .\kern.08em .\kern.08em .\kern.08em .\kern.08em .}\par} & 1818 Jan.\TD & --- & 331\\
\parbox[b]{.64\linewidth}{\raggedright \textit{Notice of the Preceding Treatise} ~\hfill\mbox{.\kern.08em .\kern.08em .\kern.08em .\kern.08em .\kern.08em .\kern.08em .\kern.08em .\kern.08em .\kern.08em .\kern.08em .\kern.08em .\kern.08em .\kern.08em .\kern.08em .}\par} & 1818 Febr.\TD & --- & 357\\
\multicolumn{4}{@{}l@{}}{\textit{Manuscript Estate.}}\\*
\multicolumn{4}{@{}p{.98\linewidth}@{}}{\TStart On the Origin and General Properties of Arithmetic–Geometric Mean Numbers.}\\
\parbox[b]{.64\linewidth}{\raggedright Art.\ 1\textemdash 8.~\hfill\mbox{.\kern.08em .\kern.08em .\kern.08em .\kern.08em .\kern.08em .\kern.08em .\kern.08em .\kern.08em .\kern.08em .\kern.08em .\kern.08em .\kern.08em .\kern.08em .\kern.08em .\kern.08em .\kern.08em .\kern.08em .\kern.08em .\kern.08em .\kern.08em .\kern.08em .\kern.08em .\kern.08em .\kern.08em .\kern.08em .\kern.08em .\kern.08em .\kern.08em .}\par} & & --- & 361\\
\end{longtable}
\endgroup

\SourceLeaf{508}


\begingroup\small
\setlength{\tabcolsep}{0pt}
\setlength{\LTleft}{0pt}\setlength{\LTright}{0pt}
\begin{longtable}{@{}l@{\hspace{.015\linewidth}}p{.075\linewidth}@{\hspace{.005\linewidth}}p{.05\linewidth}@{}}
\parbox[b]{.855\linewidth}{\raggedright On the Transcendental Functions That Arise from Differentiating Arithmetic–Geometric Means. Art.\ 9\textemdash 11.~\hfill\mbox{.\kern.08em .\kern.08em .\kern.08em .\kern.08em .\kern.08em .\kern.08em .\kern.08em .\kern.08em .\kern.08em .\kern.08em .\kern.08em .\kern.08em .\kern.08em .\kern.08em .\kern.08em .\kern.08em .\kern.08em .\kern.08em .\kern.08em .\kern.08em .\kern.08em .}\par} & Page & 372\\
\parbox[b]{.855\linewidth}{\raggedright Continuation of the Investigations of the Arithmetic–Geometric Mean. Art.\ 12\textemdash 26.~\hfill\mbox{.}\par} & --- & 375\\
\parbox[b]{.855\linewidth}{\raggedright Table. Arithmetic–Geometric Means between Unity and the Sines of Each Half Degree ~\hfill\mbox{.\kern.08em .}\par} & --- & 403\\
\parbox[b]{.855\linewidth}{\raggedright Lemniscatic Functions (I.) Represented by Power Series ~\hfill\mbox{.\kern.08em .\kern.08em .\kern.08em .\kern.08em .\kern.08em .\kern.08em .\kern.08em .\kern.08em .\kern.08em .\kern.08em .}\par} & --- & 404\\
\parbox[b]{.855\linewidth}{\raggedright Lemniscatic Functions (II.) Represented by Infinite Products and Trigonometric Series ~\hfill\mbox{.\kern.08em .\kern.08em .\kern.08em .\kern.08em .\kern.08em .\kern.08em .\kern.08em .\kern.08em .\kern.08em .\kern.08em .\kern.08em .\kern.08em .\kern.08em .\kern.08em .\kern.08em .\kern.08em .\kern.08em .\kern.08em .\kern.08em .\kern.08em .\kern.08em .\kern.08em .\kern.08em .\kern.08em .\kern.08em .\kern.08em .}\par} & --- & 413\\
\multicolumn{3}{@{}l@{}}{Concerning the Theory of the New Transcendental Functions}\\*
\parbox[b]{.855\linewidth}{\raggedright I.\quad Beginnings of the the Investigations of the New Transcendental Functions ~\hfill\mbox{.\kern.08em .\kern.08em .\kern.08em .\kern.08em .\kern.08em .\kern.08em .\kern.08em .}\par} & --- & 433\\
\parbox[b]{.855\linewidth}{\raggedright II.\quad Concerning the Theory of Transcendental Functions ~\hfill\mbox{.\kern.08em .\kern.08em .\kern.08em .\kern.08em .\kern.08em .\kern.08em .\kern.08em .\kern.08em .\kern.08em .\kern.08em .\kern.08em .\kern.08em .}\par} & --- & 436\\
\parbox[b]{.855\linewidth}{\raggedright III.\quad New Arrangement of the Theorems Concerning Those Series and Infinite Products That Belong to the Theory of Arithmetic–Geometric Means ~\hfill\mbox{.\kern.08em .\kern.08em .\kern.08em .}\par} & --- & 446\\
\parbox[b]{.855\linewidth}{\raggedright IV.\quad One Hundred Theorems on the New Transcendental Functions ~\hfill\mbox{.\kern.08em .\kern.08em .\kern.08em .\kern.08em .\kern.08em .\kern.08em .\kern.08em .\kern.08em .\kern.08em .\kern.08em .\kern.08em .\kern.08em .}\par} & --- & 461\\
\parbox[b]{.855\linewidth}{\raggedright V.\quad General Theorems ~\hfill\mbox{.\kern.08em .\kern.08em .\kern.08em .\kern.08em .\kern.08em .\kern.08em .\kern.08em .\kern.08em .\kern.08em .\kern.08em .\kern.08em .\kern.08em .\kern.08em .\kern.08em .\kern.08em .\kern.08em .\kern.08em .\kern.08em .\kern.08em .\kern.08em .\kern.08em .\kern.08em .\kern.08em .}\par} & --- & 470\\
\parbox[b]{.855\linewidth}{\raggedright The Wonderful Pentagram ~\hfill\mbox{.\kern.08em .\kern.08em .\kern.08em .\kern.08em .\kern.08em .\kern.08em .\kern.08em .\kern.08em .\kern.08em .\kern.08em .\kern.08em .\kern.08em .\kern.08em .\kern.08em .\kern.08em .\kern.08em .\kern.08em .\kern.08em .\kern.08em .\kern.08em .\kern.08em .\kern.08em .\kern.08em .\kern.08em .}\par} & --- & 481\\
\parbox[b]{.855\linewidth}{\raggedright \textit{Remarks} ~\hfill\mbox{.\kern.08em .\kern.08em .\kern.08em .\kern.08em .\kern.08em .\kern.08em .\kern.08em .\kern.08em .\kern.08em .\kern.08em .\kern.08em .\kern.08em .\kern.08em .\kern.08em .\kern.08em .\kern.08em .\kern.08em .\kern.08em .\kern.08em .\kern.08em .\kern.08em .\kern.08em .\kern.08em .\kern.08em .\kern.08em .\kern.08em .\kern.08em .\kern.08em .\kern.08em .}\par} & --- & 491\\[.7em]
\multicolumn{3}{@{}l@{}}{\textit{Two Lithographic Plates for Pages 30 and 102.}}\\*
\end{longtable}
\endgroup
\begin{center}\rule{7em}{0.4pt}\end{center}

\clearpage\fancyhead[L]{}\pdfbookmark[0]{Printer’s Imprint}{original-colophon}\SourceLeaf{509}

\vspace*{0.24\textheight}
\begin{center}
{\large GÖTTINGEN,}\par
\medskip
\textbf{PRINTED AT THE DIETERICH\\UNIVERSITY PRINTING HOUSE}\par
\smallskip
\textsc{W. Fr. Kaestner.}
\end{center}
\clearpage
\endgroup
